% Part: second-order-logic
% Chapter: sol-and-sets
% Section: introduction

\documentclass[../../../include/open-logic-section]{subfiles}

\begin{document}

\olfileid[pt]{sol}{set}{int}

\olsection{Introdução}

Como a lógica de segunda ordem pode quantificar tanto sobre subconjuntos do
domínio quanto sobre funções, é de esperar que pelo menos alguma teoria dos
conjuntos possa ser desenvolvida na lógica de segunda ordem. Por
``desenvolvida'', queremos dizer que é possível expressar propriedades e
afirmações da teoria dos conjuntos na lógica de segunda ordem, e que isso é
possível sem nenhum vocabulário não lógico especial para conjuntos (p.ex., o !!{predicate} de pertinência da teoria dos conjuntos). Por
exemplo, podemos definir uniões e interseções de conjuntos e a relação de
inclusão, mas também comparar os tamanhos de conjuntos e enunciar resultados
como o teorema de Cantor.

\end{document}

